% VIGTIGT! For at dokumentet kan compile, skal man benytte XeLaTeX eller LuaLaTeX som compiler. Dette kan gøres i Overleaf ved at gå til Menu -> Settings -> Compiler -> Vælg XeLaTeX/LuaLaTeX
\documentclass[t,24pt,aspectratio=169]{beamer}

\usepackage{KUstyle}
\usepackage{multicol} % Pakke til multiple kolonner

\toplinje{Tekst på toplinjen} % Teksten på toplinjen. Hvis ingen tekst ønskes, skal kommandoen fjernes

\begin{document}

% Det første slide. Her kan man f.eks. ændre hovedtitel, undertitlen, oplægsholder, KU-enhed, dato og baggrundsbilledet
{
\setbeamertemplate{background}{\includegraphics[width=\paperwidth,height=\paperheight]{billeder/forside.jpg}}
\begin{frame}
    \begin{textblock*}{\textwidth}(0.57\textwidth,0.1\textheight)
        \begin{beamercolorbox}[wd=7.8cm,ht=7.3cm,sep=0.5cm]{hvidbox}
            \fontsize{5}{10}\fontfamily{ptm}\selectfont \textls[200]{KØBENHAVNS UNIVERSITET}
            \noindent\textcolor{KUrod}{\rule{6.8cm}{0.4pt}}
        \end{beamercolorbox}
    \end{textblock*}
    \begin{textblock*}{\textwidth}(0.57\textwidth,0.1\textheight)
        \begin{beamercolorbox}[wd=7.8cm,sep=0.5cm]{hvidbox}
                \Huge \textcolor{KUrod}{Hovedtitel}
                \vspace{0.5cm}
                \par
                \Large Dette er undertitlen
                \vspace{0.5cm}
                \par
                \normalsize Oplægsholder, KU-enhed, dato
        \end{beamercolorbox}
    \end{textblock*}
    \begin{textblock}{1}(14.5,11.38)
        \includegraphics[width=1cm]{KU/KU-logo.png}
    \end{textblock}
\end{frame}
}

% Et standard slide. Det er vigtigt at enten [hvid], [rod] eller [billede] er inkluderet efter \begin{frame}. Hvis man benytter [billede] er det også vigtigt at bruge \setbeamertemplate{background} kommandoen til at angive baggrundsbilledet (længere nede kan det ses, hvordan det helt præcist gøres)
\begin{frame}[hvid]
\frametitle{Hvid standard slide}
Lorem ipsum dolor sit amet, consectetur adipiscing elit. Quisque viverra rhoncus urna, et sodales magna efficitur et. Integer tempor augue eget ultricies malesuada. Praesent maximus sapien et tellus venenatis, nec consequat sem dictum. In lectus tellus, sagittis at massa in, tempor pretium quam. Pellentesque vel scelerisque mi.
\\~\\
Sed vel ligula egestas, commodo lacus in, cursus velit. Nam faucibus felis et nibh vehicula, eget aliquam justo facilisis. Aenean a dui sed libero malesuada dictum sit amet quis turpis. Nulla consectetur lorem at augue lacinia, ut rutrum neque vulputate. Aliquam ultrices suscipit felis vel pellentesque. Donec lacinia orci nec feugiat rhoncus. Nunc congue at velit vitae posuere. In hac habitasse platea dictumst.
\end{frame}

% Det efterfølgende er forskellige eksempler på slides
\begin{frame}[hvid]
\frametitle{Liste med punkter}
    \begin{itemize}
        \item Punkt 1
        \item Punkt 2
            \begin{itemize}
                \item Punkt 2.1
                \item Punkt 2.2
            \end{itemize}
        \item Punkt 3
    \end{itemize}
\end{frame}

\begin{frame}[hvid]
\frametitle{Slide med 2 kolonner}
    \begin{multicols}{2}
        \begin{itemize}
            \item Punkt 1
            \item Punkt 2
        \end{itemize}
        \vfill\null
        \columnbreak
        \begin{itemize}
            \item Punkt 3
            \item Punkt 4
            \item Punkt 5
        \end{itemize}
    \end{multicols}
\end{frame}

\begin{frame}[hvid]
\frametitle{Slide med en kolonne og et billede}
    \begin{multicols}{2}
        \begin{itemize}
            \item Punkt 1
            \item Punkt 2
        \end{itemize}
        \vfill\null
        \columnbreak
            \includegraphics[width=6.6cm,height=5.9cm,keepaspectratio]{billeder/1.jpg}
    \end{multicols}
\end{frame}

\begin{frame}[hvid]
\frametitle{Slide med stort billede}
    \vspace{0.5cm}
    \includegraphics[width=14.5cm,height=6cm,keepaspectratio]{billeder/2.jpg}
\end{frame}

% Her er et eksempel med [billede], hvori man angiver et baggrundsbillede
{
\setbeamertemplate{background}{\includegraphics[width=\paperwidth,height=\paperheight]{billeder/3.jpg}}
\begin{frame}[billede]
    \begin{textblock*}{\textwidth}(0\textwidth,0.77\textheight)
        \begin{beamercolorbox}[wd=7.8cm,sep=0.3cm]{rodbox}
        Et lille mærke med en vilkårlig tekst
        \end{beamercolorbox}
    \end{textblock*}
\end{frame}
}

{
\setbeamertemplate{background}{\includegraphics[width=\paperwidth,height=\paperheight]{billeder/4.jpg}}
\begin{frame}[billede]
    \begin{textblock*}{\textwidth}(0.57\textwidth,0.4\textheight)
        \begin{beamercolorbox}[wd=7.8cm,sep=0.3cm]{rodbox}
        Dette er et stort mærke. Her kan man f.eks. have en liste
            \begin{itemize}
                \item Punkt 1
                \item Punkt 2
                \item Punkt 3
                \item Punkt 4
            \end{itemize}
        \end{beamercolorbox}
    \end{textblock*}
\end{frame}
}

\begin{frame}[rod]
    \fontsize{80}{0}\selectfont ``
    \\
    \vspace{-0.6cm}
    \fontsize{32}{0}\selectfont Prediction is very difficult, especially if it's about the future.
    \\
    \vspace{2.8cm}
    \fontsize{18}{0}\selectfont Niels Bohr
\end{frame}

\begin{frame}[rod]
    \vspace{0.9cm}
    \resizebox{\textwidth}{!}{
        \begin{tabular}{r l}
            \textbf{Research} & Forskning \\
            \textbf{Education} & Uddannelse \\
            \textbf{Exchange of knowdledge} & Forskningsformidling \\
        \end{tabular}
    }
\end{frame}

\end{document}